\section{\texorpdfstring{Sustancias, interacciones y cambios visuales en papel: introducción al diseño factorial \(2^2\)}{Sustancias, interacciones y cambios visuales en papel: introducción al diseño factorial 2 por 2}}

\subsection{Introducción}

Las organizaciones frecuentemente deben tomar decisiones acerca de qué materiales, procesos o tratamientos producen mejores resultados. Sin embargo, cuando intervienen varios factores simultáneamente, resulta difícil determinar cuál de ellos es responsable de los cambios observados.

Una estrategia ampliamente utilizada para abordar este problema consiste en emplear diseños experimentales factoriales. Estos diseños permiten estudiar simultáneamente varios factores, evaluar sus efectos individuales y analizar posibles interacciones entre ellos.

En esta práctica se utilizarán soluciones elaboradas con cúrcuma y cloro comercial aplicadas sobre papel adhesivo (\textit{post-it}) para generar cambios visuales observables. Aunque el sistema experimental es sencillo, reproduce la lógica utilizada en numerosos experimentos industriales, agrícolas y científicos.

El objetivo principal no consiste en estudiar las propiedades químicas de la cúrcuma o del cloro, sino comprender cómo se construye un diseño factorial, cómo se establecen tratamientos experimentales, cuál es la función de un blanco experimental y cómo se registran e interpretan los resultados obtenidos.

Además, se incorporará el uso de un luxómetro para transformar observaciones visuales en mediciones cuantitativas, permitiendo analizar los tratamientos mediante herramientas estadísticas objetivas.

\subsection{Objetivo general}

Comprender la estructura y aplicación de un diseño factorial $2^2$ mediante el análisis de cambios visuales y de transmisión luminosa producidos por soluciones de cúrcuma y cloro comercial sobre papel.

\subsection{Objetivos particulares}

\begin{itemize}
\item Identificar factores, niveles y tratamientos experimentales.
\item Comprender la función de un blanco experimental.
\item Analizar la importancia de las repeticiones experimentales.
\item Registrar y comparar cambios visuales producidos por diferentes tratamientos.
\item Introducir el concepto de interacción entre factores.
\item Medir cuantitativamente la transmisión luminosa utilizando un luxómetro.
\item Generar evidencia experimental para la toma de decisiones.
\end{itemize}

\subsection{Fundamento teórico}

Los diseños factoriales permiten estudiar simultáneamente el efecto de varios factores sobre una o más variables respuesta. Cuando se analizan dos factores y cada uno posee dos niveles, se obtiene un diseño factorial:

\[
2^2
\]

Este diseño genera cuatro combinaciones experimentales diferentes y permite estimar:

\begin{itemize}
\item el efecto principal del primer factor,
\item el efecto principal del segundo factor,
\item y la posible interacción entre ambos factores.
\end{itemize}

La principal ventaja de este enfoque es que permite obtener más información utilizando un número relativamente pequeño de tratamientos experimentales.

\subsection{Contexto del problema}

Una empresa dedicada a la fabricación de materiales decorativos desea desarrollar papeles con distintos efectos visuales. Antes de invertir en producción necesita determinar qué combinación de sustancias genera los cambios más interesantes y qué tratamientos modifican la transmisión de luz a través del papel.

Para responder estas preguntas se diseñará un experimento controlado utilizando cúrcuma y cloro comercial como factores experimentales.

\subsection{Factores experimentales}

En esta práctica se estudiarán dos factores.

\subsubsection{Factor A: concentración de cúrcuma}

\begin{table}[H]
\centering
\begin{tabular}{cc}
\hline
Nivel & Concentración \\
\hline
Bajo (-) & 1 g \\
Alto (+) & 4 g \\
\hline
\end{tabular}
\caption{Niveles del factor cúrcuma.}
\end{table}

\subsubsection{Factor B: concentración de cloro comercial}

\begin{table}[H]
\centering
\begin{tabular}{cc}
\hline
Nivel & Concentración \\
\hline
Bajo (-) & 5 g \\
Alto (+) & 10 g \\
\hline
\end{tabular}
\caption{Niveles del factor cloro comercial.}
\end{table}

\subsection{Variables controladas}

Con el propósito de que los cambios observados puedan atribuirse a los factores estudiados, todos los tratamientos mantendrán constantes las siguientes condiciones:

\begin{itemize}
\item 350 g de agua potable.
\item 20 g de sal.
\item Tres aspersiones por unidad experimental.
\item Distancia de aplicación de 10 cm.
\item Mismo tipo de papel adhesivo.
\item Mismo tiempo de secado.
\item Mismo tamaño de la pieza de madera utilizada para cubrir parcialmente el papel.
\end{itemize}

\subsection{Diseño experimental}

La estructura factorial utilizada se muestra en la Tabla \ref{tab:factorial}.

\begin{table}[H]
\centering
\caption{Tratamientos factoriales del diseño $2^2$.}
\label{tab:factorial}
\begin{tabular}{ccc}
\hline
Tratamiento & Cúrcuma & Cloro \\
\hline
T1 & 1 g & 5 g \\
T2 & 1 g & 10 g \\
T3 & 4 g & 5 g \\
T4 & 4 g & 10 g \\
\hline
\end{tabular}
\end{table}

\subsection{Blanco experimental}

Además de los tratamientos factoriales se incorporará un blanco experimental compuesto únicamente por:

\[
350 \text{ g de agua potable}
\]

y

\[
20 \text{ g de sal}
\]

sin cúrcuma y sin cloro comercial.

La función del blanco es proporcionar una referencia para determinar si los cambios observados son consecuencia de los factores estudiados o simplemente resultado de la humedad o del proceso de aplicación.

\subsection{Unidad experimental}

La unidad experimental corresponderá a un post-it individual.

Cada tratamiento contará con seis repeticiones independientes.

Por lo tanto, el número total de unidades experimentales será:

\[
5 \times 6 = 30
\]

donde los cinco tratamientos corresponden a los cuatro tratamientos factoriales y al blanco experimental.

\subsection{Preparación de las soluciones}

Las soluciones se prepararán al menos 24 horas antes de la aplicación experimental.

El objetivo de este periodo de reposo es permitir una adecuada dispersión de los componentes y garantizar condiciones homogéneas entre tratamientos.

Cada solución contendrá:

\begin{itemize}
\item 350 g de agua potable,
\item 20 g de sal,
\item y las cantidades correspondientes de cúrcuma y cloro según el tratamiento asignado.
\end{itemize}

\begin{table}[H]
\centering
\begin{tabular}{cccc}
\hline
Tratamiento & Cúrcuma & Cloro & Sal \\
\hline
T1 & 1 g & 5 g & 20 g \\
T2 & 1 g & 10 g & 20 g \\
T3 & 4 g & 5 g & 20 g \\
T4 & 4 g & 10 g & 20 g \\
B1 & 0 g & 0 g & 20 g \\
\hline
\end{tabular}
\caption{Composición de los tratamientos experimentales.}
\end{table}

\subsection{Procedimiento experimental}

\begin{enumerate}
\item Colocar un post-it sobre una superficie plana.

\item Cubrir aproximadamente la mitad de la superficie mediante una pieza de madera.

\item Colocar el atomizador inclinado respecto al papel.

\item Mantener una distancia aproximada de:

\[
10 \text{ cm}
\]

entre el atomizador y la superficie experimental.

\item Aplicar tres aspersiones sobre la muestra.

\item Retirar cuidadosamente la pieza de madera.

\item Dejar secar la muestra bajo las mismas condiciones ambientales.

\item Registrar fotografías y observaciones.
\end{enumerate}

\subsection{Variables respuesta}

La práctica considera variables tanto cualitativas como cuantitativas.

\subsubsection{Variables cualitativas}

\begin{itemize}
\item Intensidad visual del color.
\item Uniformidad de la coloración.
\item Contraste entre zonas protegidas y expuestas.
\item Atractivo visual.
\end{itemize}

\subsubsection{Variables cuantitativas}

La variable principal será la transmisión luminosa medida mediante un luxómetro.

\[
L=\text{lux medidos}
\]

Esta variable permitirá cuantificar la cantidad de luz que atraviesa cada unidad experimental.

\subsection{Medición mediante luxómetro}

Una vez que las muestras hayan secado completamente, se procederá a medir la transmisión luminosa utilizando un luxómetro.

Para ello se utilizará una fuente de luz constante colocada debajo de la muestra.

El montaje experimental será:

\[
\text{Fuente de luz} \rightarrow \text{Post-it} \rightarrow \text{Luxómetro}
\]

Primero se registrará la intensidad luminosa sin muestra:

\[
L_0
\]

Posteriormente se registrará la intensidad luminosa atravesando cada post-it tratado:

\[
L_i
\]

Con estos valores podrá calcularse el porcentaje de bloqueo de luz:

\[
\%B=
\frac{L_0-L_i}{L_0}\times100
\]

donde:

\begin{itemize}
\item $L_0$ corresponde a la lectura sin muestra.
\item $L_i$ corresponde a la lectura con la muestra.
\end{itemize}

\subsection{Tabla de registro experimental}

\begin{table}[H]
\centering
\resizebox{0.98\linewidth}{!}{%
\begin{tabular}{ccccccccc}
\hline
Trat. & Rep. & Lux & \% Bloqueo & Color & Contraste & Uniformidad & Atractivo & Obs. \\
\hline
T1 & 1 & & & & & & & \\
T1 & 2 & & & & & & & \\
T1 & 3 & & & & & & & \\
T1 & 4 & & & & & & & \\
T1 & 5 & & & & & & & \\
T1 & 6 & & & & & & & \\
\hline
\end{tabular}%
}
\caption{Formato de captura de datos experimentales.}
\end{table}

\subsection{Análisis de resultados}

Con los datos obtenidos se podrá:

\begin{itemize}
\item Calcular promedios por tratamiento.
\item Comparar tratamientos utilizando diagramas de caja.
\item Construir gráficas de efectos principales.
\item Analizar posibles interacciones entre factores.
\item Realizar un ANOVA factorial utilizando la variable lux como respuesta principal.
\end{itemize}

\subsubsection{\texorpdfstring{Análisis estadístico mediante ANOVA factorial \(2^2\)}{Análisis estadístico mediante ANOVA factorial 2 por 2}}

Para analizar cuantitativamente los resultados se utiliza un modelo estadístico lineal que permite descomponer la variabilidad total observada en componentes atribuibles a cada factor y a su interacción.

\paragraph{Modelo estadístico}

\[
Y_{ijk} = \mu + A_i + B_j + (AB)_{ij} + \varepsilon_{ijk}
\]

donde:
\begin{itemize}
\item $Y_{ijk}$: variable respuesta (porcentaje de bloqueo de luz) para la réplica $k$ en la combinación de niveles $i$ de cúrcuma y $j$ de cloro.
\item $\mu$: media global del experimento.
\item $A_i$: efecto del factor cúrcuma en el nivel $i$ (bajo o alto).
\item $B_j$: efecto del factor cloro en el nivel $j$ (bajo o alto).
\item $(AB)_{ij}$: efecto de interacción entre cúrcuma y cloro.
\item $\varepsilon_{ijk}$: error experimental, que se asume aleatorio, independiente y distribuido normalmente con media cero y varianza constante.
\end{itemize}

\paragraph{Codificación de niveles y tratamientos}

Para facilitar los cálculos, se asigna el valor $-$ al nivel bajo y $+$ al nivel alto de cada factor.

\begin{table}[H]
\centering
\resizebox{0.98\linewidth}{!}{%
\begin{tabular}{ccccc}
\hline
Tratamiento & Cúrcuma (A) & Cloro (B) & Codificación & Promedio \% bloqueo \\
\hline
T1 & 1 g ($-$) & 5 g ($-$) & $(-,-)$ & $\overline{Y}_1$ \\
T2 & 1 g ($-$) & 10 g (+) & $(-,+)$ & $\overline{Y}_2$ \\
T3 & 4 g (+) & 5 g ($-$) & $(+,-)$ & $\overline{Y}_3$ \\
T4 & 4 g (+) & 10 g (+) & $(+,+)$ & $\overline{Y}_4$ \\
\hline
\end{tabular}%
}
\caption{Codificación de niveles y promedios por tratamiento.}
\end{table}

\paragraph{Cálculo de efectos principales e interacción}

Los efectos se estiman mediante contrastes ortogonales utilizando los signos $+$ y $-$.

\textbf{Efecto principal de A (cúrcuma):}

Compara el promedio de los tratamientos con nivel alto de cúrcuma contra los de nivel bajo.

\[
\text{Efecto}_A = \frac{(\overline{Y}_3 + \overline{Y}_4)}{2} - \frac{(\overline{Y}_1 + \overline{Y}_2)}{2}
\]

En forma de contraste con signos:

\[
\text{Efecto}_A = \frac{-\overline{Y}_1 - \overline{Y}_2 + \overline{Y}_3 + \overline{Y}_4}{2}
\]

\textbf{Efecto principal de B (cloro):}

Compara el promedio de los tratamientos con nivel alto de cloro contra los de nivel bajo.

\[
\text{Efecto}_B = \frac{(\overline{Y}_2 + \overline{Y}_4)}{2} - \frac{(\overline{Y}_1 + \overline{Y}_3)}{2}
\]

En forma de contraste con signos:

\[
\text{Efecto}_B = \frac{-\overline{Y}_1 + \overline{Y}_2 - \overline{Y}_3 + \overline{Y}_4}{2}
\]

\textbf{Interacción AB:}

Mide si el efecto de un factor depende del nivel del otro factor.

\[
\text{Efecto}_{AB} = \frac{(\overline{Y}_1 + \overline{Y}_4)}{2} - \frac{(\overline{Y}_2 + \overline{Y}_3)}{2}
\]

En forma de contraste con signos:

\[
\text{Efecto}_{AB} = \frac{+\overline{Y}_1 - \overline{Y}_2 - \overline{Y}_3 + \overline{Y}_4}{2}
\]

\paragraph{Sumas de cuadrados (SC)}

Cada contraste se construye con las sumas de las repeticiones, no con los promedios. Sea $n$ el número de repeticiones por tratamiento (en este caso $n=6$).

\textbf{Contraste para A:}

\[
\text{Contraste}_A = -\sum_{k=1}^{n} Y_{1k} -\sum_{k=1}^{n} Y_{2k} +\sum_{k=1}^{n} Y_{3k} +\sum_{k=1}^{n} Y_{4k}
\]

\textbf{Suma de cuadrados para A:}

\[
SC_A = \frac{(\text{Contraste}_A)^2}{4n}
\]

\textbf{Análogamente para B y AB:}

\[
SC_B = \frac{(\text{Contraste}_B)^2}{4n}, \quad SC_{AB} = \frac{(\text{Contraste}_{AB})^2}{4n}
\]

\textbf{Suma de cuadrados total:}

\[
SC_{\text{total}} = \sum_{i=1}^{4} \sum_{k=1}^{n} Y_{ik}^2 - \frac{(\sum_{i=1}^{4} \sum_{k=1}^{n} Y_{ik})^2}{4n}
\]

\textbf{Suma de cuadrados del error:}

\[
SC_{\text{error}} = SC_{\text{total}} - (SC_A + SC_B + SC_{AB})
\]

\paragraph{Cuadrados medios (CM) y estadístico F}

Los cuadrados medios se obtienen dividiendo cada suma de cuadrados entre sus grados de libertad (gl):

\begin{itemize}
\item Factor A: gl = 1
\item Factor B: gl = 1
\item Interacción AB: gl = 1
\item Error: gl = $4(n-1) = 4(6-1) = 20$
\item Total: gl = $4n - 1 = 23$
\end{itemize}

\[
CM_A = \frac{SC_A}{1}, \quad CM_B = \frac{SC_B}{1}, \quad CM_{AB} = \frac{SC_{AB}}{1}, \quad CM_{\text{error}} = \frac{SC_{\text{error}}}{4(n-1)}
\]

El estadístico $F$ para cada fuente se calcula como:

\[
F_A = \frac{CM_A}{CM_{\text{error}}}, \quad F_B = \frac{CM_B}{CM_{\text{error}}}, \quad F_{AB} = \frac{CM_{AB}}{CM_{\text{error}}}
\]

Cada $F$ se compara con el valor crítico de la distribución $F$ con 1 y $4(n-1)$ grados de libertad para determinar si el efecto es estadísticamente significativo.

\paragraph{Tabla de ANOVA}

\begin{table}[H]
\centering
\begin{tabular}{lccccc}
\hline
Fuente de variación & SC & gl & CM & F & p-valor \\
\hline
Cúrcuma (A) & $SC_A$ & 1 & $CM_A$ & $F_A$ & \\
Cloro (B) & $SC_B$ & 1 & $CM_B$ & $F_B$ & \\
Interacción AB & $SC_{AB}$ & 1 & $CM_{AB}$ & $F_{AB}$ & \\
Error & $SC_{\text{error}}$ & $4(n-1)$ & $CM_{\text{error}}$ & & \\
\hline
Total & $SC_{\text{total}}$ & $4n-1$ & & & \\
\hline
\end{tabular}
\caption{Tabla de ANOVA para el diseño factorial $2^2$.}
\end{table}

\paragraph{Interpretación de resultados}

\begin{itemize}
\item Si $F_A$ es mayor que el valor crítico (o el p-valor < 0.05), se concluye que la concentración de cúrcuma afecta significativamente el bloqueo de luz.
\item Si $F_B$ es significativo, el cloro tiene un efecto significativo.
\item Si $F_{AB}$ es significativo, existe interacción: el efecto de un factor depende del nivel del otro.
\end{itemize}

La ventaja de este análisis es que transforma observaciones visuales subjetivas en decisiones objetivas basadas en evidencia estadística.

\subsubsection{Ejemplo ilustrativo (datos hipotéticos)}

Supóngase que con $n=6$ repeticiones se obtienen los siguientes promedios de porcentaje de bloqueo de luz:

\begin{itemize}
\item T1 ($-$,$-$): 10\%
\item T2 ($-$,$+$): 25\%
\item T3 ($+$,$-$): 30\%
\item T4 ($+$,$+$): 60\%
\end{itemize}

Aplicando las fórmulas:

\[
\text{Efecto}_A = \frac{30+60}{2} - \frac{10+25}{2} = 45 - 17.5 = 27.5\%
\]

\[
\text{Efecto}_B = \frac{25+60}{2} - \frac{10+30}{2} = 42.5 - 20 = 22.5\%
\]

\[
\text{Efecto}_{AB} = \frac{10+60}{2} - \frac{25+30}{2} = 35 - 27.5 = 7.5\%
\]

Interpretación:
\begin{itemize}
\item Pasar de 1 g a 4 g de cúrcuma aumenta el bloqueo en 27.5\%.
\item Pasar de 5 g a 10 g de cloro aumenta el bloqueo en 22.5\%.
\item La interacción positiva (7.5\%) indica que el efecto combinado de ambos factores es ligeramente mayor que la suma de sus efectos individuales.
\end{itemize}

Con un ANOVA se determinaría si estos efectos son estadísticamente significativos.

\subsection{Preguntas de análisis}

\begin{enumerate}
\item ¿Qué diferencia existe entre un factor y un nivel?
\item ¿Por qué este experimento corresponde a un diseño factorial $2^2$?
\item ¿Qué función cumple el blanco experimental?
\item ¿Por qué se utilizan repeticiones?
\item ¿Qué tratamiento produjo el mayor cambio visual?
\item ¿Qué tratamiento produjo el mayor bloqueo de luz?
\item ¿Existe evidencia de interacción entre la cúrcuma y el cloro?
\item ¿Qué ocurriría si se eliminara el blanco experimental?
\item ¿Qué ventaja tiene utilizar un luxómetro en lugar de únicamente observaciones visuales?
\item ¿Qué variable utilizarías para realizar un ANOVA y por qué?
\end{enumerate}

\subsection{Reflexión final}

Aunque el sistema experimental utiliza materiales sencillos, reproduce los principios fundamentales empleados en investigación científica y desarrollo tecnológico. Comprender cómo se diseñan los tratamientos, cómo se controlan las variables y cómo se interpretan los resultados constituye una habilidad esencial para la toma de decisiones basada en evidencia.

La incorporación del luxómetro permite además distinguir entre observaciones subjetivas y mediciones instrumentales, fortaleciendo la calidad de las conclusiones obtenidas. El uso del ANOVA factorial con la variable cuantitativa de porcentaje de bloqueo de luz permite decidir estadísticamente si los efectos observados son reales o debidos al azar, tal como se hace en la industria farmacéutica, agrícola y de manufactura para optimizar procesos.